\scanpage{16}

\deleted{Or toutes ces assertions découlent facilement de la suivante : tt diagramme \inserted{$\delta$} de type $\Theta\in d''$ dans $D$ provient d'un diagramme $\delta_j$ dans un des $C_j$ $(j<i)$, \uline{et} que l'ens. d'indices des $j<i$ est grand devant $\mathrm{card}\,\mathrm{Fl}\,\Theta$ : on voit alors que la $\varinjlim$ de $\unsure{\overrightarrow{\delta_j}|\delta_j|D}$ dans $C_j$ donne une limite dans $D$ [car les $\varprojlim$ de type $\Theta$ commutent aux $\varinjlim$ grandes devant \inserted{$\mathrm{card}$} $\mathrm{Fl}\,\Theta$)].}

\deleted{b) $\alpha_i:C\unsure{-|\to} C_i$ est $\delta'$-exact, i.e. si $(x\overset{p_\lambda}{-}x_\lambda)$ est dans $\delta'$, alors $(\alpha_i(x)\overset{\alpha_i(p_\lambda)}{-}\alpha_i(x_\lambda))$ est un cône proj. \uline{exact}.}

% slanted
\margin{\uline{\unsure{redonner|redonne}}}

\deleted{Or cela résulte aussitôt de la description des flèches de $\varinjlim_{j<i\,(\mathrm{Cat})} C_i$, de a) et b), et du fait que $i=\{j\mid j<i\}$ est grand devant $\mathrm{card}\,\mathrm{Fl}\,\Theta$ -- de sorte que \unsure{tte|tt} syst \inserted{``compatible''} de flèches $z-\alpha_i(x_\lambda)$ provient d'un syst. analogue dans un $C_j$, $j<i$ \inserted{satisfaisant les compatibilités analogues}, et tte double flèche $u,v:z-\alpha_i(u)$ provient d'une double flèche analogue dans un $C_j$, l'égalité $\alpha_i(p_\lambda)u=\unsure{d_i|\alpha_i}(p_\lambda)v$ pour tt $\lambda$ étant déjà vérifiée.}

\deleted{Cela achève la démonstration.}

% slanted
\margin{\uline{Ce ci résoud à la fois le \deleted{\unsure{seul}} Pb 1 et le Pb 2}}

\deleted{\uline{Th} Soient $d',d''$ des petits ensembles de diagrammes, \deleted{\unsure{et}} $\Delta=\Delta_{(d',d'')}$, $C$ une petite catégorie \inserted{munie d'un ens $(\delta',\delta'')$ de cônes projectifs et inductifs}. Alors \deleted{\unsure{le}} il existe une $\Delta$-enveloppe de $(C,\delta',\delta'')$, \deleted{et celle-ci est ens$^{\mathrm{t}}$ petite.} i.e. un foncteur $u:C-C'$ avec $C'\in\mathrm{Ob}\,\Delta$, $i$ $(\delta',\delta'')$ exact, et universel pour ces propriétés [i.e. $\forall E\in\mathrm{Ob}\,\Delta$, $\varphi\mapsto\varphi\circ i$ induit une équivalence \deleted{$\underline{\mathrm{Hom}}_{\Delta,(\delta',\delta'')}(C$} $\underline{\mathrm{Hom}}_{\Delta,(i(\delta'),i(\delta''))}(C',E)\xrightarrow{\sim}\underline{\mathrm{Hom}}_{\unsure{\Delta,}(\delta',\delta'')}(C,E)$.}

\deleted{\unsure{Soit} Soit $\rho$ le foncteur $\Delta_{(\varnothing,d'')}$-enveloppe,\\ $\sigma$ \hspace{.7em}---\hspace{.7em} $\Delta_{(d',\varnothing)}$-enveloppe.}

\deleted{On construit une suite transfinie de catégories $C_i$ sous $C$, chacune munie d'un ens \inserted{$\delta_i=$} $(\delta'_i,\delta''_i)$ de cônes proj. et inductifs. On pose}

\[
\begin{aligned}
&\deleted{C_0=C,\quad\deleted{\delta'_0=\delta''_0=\varnothing}\quad(\delta'_0,\delta''_0)=(\delta',\delta'')}\\
&\deleted{(C_{i+1},\inserted{\delta_{i+1})}\text{ déduit de }(C_i,\delta_i)\text{ ainsi :}}
\end{aligned}
\]

% archive sheet number: 15
